%%%%%%%%%%%%%%%%%%%%%%%%%%%%%%%%%%%%%%%%%%%%%%%%%%%%%%%%%%%%%%%%%%%%%%%%%%%%%%%%
% ScaleNav -- ICRA-style manuscript (revised draft)
% Compile with ./build.sh
% Figures under pics/ are rendered when present; otherwise a labeled
% placeholder box is shown so the manuscript compiles standalone.
%%%%%%%%%%%%%%%%%%%%%%%%%%%%%%%%%%%%%%%%%%%%%%%%%%%%%%%%%%%%%%%%%%%%%%%%%%%%%%%%

\documentclass[letterpaper,10pt,conference]{ieeeconf}
\IEEEoverridecommandlockouts
\overrideIEEEmargins

\usepackage{amsmath}
\usepackage{amssymb}
\usepackage{booktabs}
\usepackage{graphicx}
\usepackage{xcolor}
\usepackage{placeins}
\usepackage{xstring}

\newcommand{\vect}[1]{\boldsymbol{#1}}
\newcommand{\mat}[1]{\mathbf{#1}}
\newcommand{\R}{\mathbb{R}}
\newcommand{\blank}{--}

% Render the figure when available; otherwise insert a labeled placeholder so
% that the manuscript compiles without the pics/ directory.
\newcommand{\safefig}[2][]{%
  \IfFileExists{#2}{\includegraphics[#1]{#2}}%
  {\begingroup\setlength{\fboxsep}{4pt}%
   \framebox[\dimexpr\linewidth-2\fboxsep-2\fboxrule]{%
     \parbox{0.88\linewidth}{\centering\vspace{1.6em}%
       \footnotesize\textit{[Figure placeholder]}\\[2pt]%
       \begingroup\footnotesize\ttfamily
       \edef\safefigpath{\detokenize{#2}}%
       \StrSubstitute{\safefigpath}{/}{/\allowbreak}[\safefigpathA]%
       \StrSubstitute{\safefigpathA}{_}{\_\allowbreak}[\safefigpathB]%
       \safefigpathB\endgroup%
       \vspace{1.6em}}}\endgroup}}

\title{\LARGE \bf
ScaleNav: A Plug-and-Play Topological Route Layer\\
for Language-Guided Aerial Navigation
}

\author{First Author$^{1}$ and Second Author$^{2}$%
\thanks{This work was not supported by any organization.}%
\thanks{$^{1}$First Author is with the Affiliation,
        {\tt\small first.author@example.com}}%
\thanks{$^{2}$Second Author is with the Affiliation,
        {\tt\small second.author@example.com}}%
}

\IEEEaftertitletext{%
  \begin{center}
  \safefig[width=0.98\textwidth]{pics/teaser.pdf}
  \vspace{3pt}
  \parbox{0.98\textwidth}{\footnotesize
  \textbf{Fig. 1.}~ScaleNav in Map2 at synchronized evidence time $t^*$.
  (a)~RGB, semantic score, and depth show semantic response where
  depth is clipped beyond $20$~m. (b)~A shallow aerial view of the mission over
  surface-sampled UE mesh truth. Graph-node fill encodes each node's semantic
  score; crosses mark far-field semantic-node
  positions that lie within the planner's three-dimensional influence radius
  of this graph slice. The persistent graph is shown in a $\pm2$~m
  navigation-height slice. The
  graph/path state is within $90$~ms of $t^*$.
  \texttt{mission\_goal}, \texttt{frontier\_goal}, and \texttt{local\_goal} denote the task,
  route, and execution layers. Semantic risk changes verified-edge costs;
  crosses denote provisional frontiers beyond the measured backbone.}
  \end{center}%
  \setcounter{figure}{1}%
}

\begin{document}
\maketitle
\thispagestyle{empty}
\pagestyle{empty}

%%%%%%%%%%%%%%%%%%%%%%%%%%%%%%%%%%%%%%%%%%%%%%%%%%%%%%%%%%%%%%%%%%%%%%%%%%%%%%%%
\begin{abstract}
Goal-conditioned local planners react quickly to observed geometry but do not
retain the route memory needed to circumnavigate obstacles larger than their
sensing horizon.  ScaleNav is a plug-and-play topological route layer that
adds this memory without retraining or changing the core algorithm of the
local planner.  From depth and odometry, it maintains persistent verified
free-space bubbles and collision-checked witness paths while restricting each
online search to a local working set.  Open-vocabulary image responses add
provisional far-field frontiers and continuous language-conditioned cost along
verified witnesses.  A topology-graph GCN learns long-range frontier-direction
selection from the online skeleton, while A* and geometric checks remain
authoritative.  The accepted witness is exposed as a
moving world-frame local goal: YOPO consumes it directly, while thin ROS
adapters relay the same interface to EGO-Planner and SUPER.  No
route-conditioned model or planner retraining is required.  In Map2, the
standalone YOPO, EGO-Planner, and SUPER batches each completed 0/10
missions, whereas their ScaleNav-guided variants completed 10/10, 7/10, and
8/10, respectively.  Across the persistent-map profile, the local graph and
active A* workloads remain 214/249 and 185/215 nodes at P50/P95.  The result
is an architectural separation: topology remembers and selects the route,
while an interchangeable local planner executes only the next goal.
\end{abstract}

%%%%%%%%%%%%%%%%%%%%%%%%%%%%%%%%%%%%%%%%%%%%%%%%%%%%%%%%%%%%%%%%%%%%%%%%%%%%%%%%
\section{INTRODUCTION}
Flying quickly through an unknown environment forces a robot to reconcile
processes that naturally operate at very different rates.  A local policy must
convert the current depth image and vehicle state into a dynamically feasible
command within milliseconds
\cite{rappids,nanomap,egoplanner,yopo,loquercio}, whereas building a reusable
map, reasoning over long-horizon alternatives, and running open-vocabulary
perception all take orders of magnitude longer
\cite{voxblox,clip}.  Learned one-stage planners
\cite{yopo,loquercio} have shown
how much a fast reactive layer alone can achieve; yet a single observation
cannot preserve an early left--right commitment around a long obstacle, and
once the fork disappears from view a memoryless policy is free to undo it.
Topological memories in ground navigation have repeatedly demonstrated that
retaining route structure, rather than reconstructing it from scratch at every
tick, is what converts local reactivity into reliable long-range behavior
\cite{sptm,ving}.

A complementary opportunity has emerged on the perception side.  RGB carries
evidence that depth does not: a power line, a glass facade, or a distant
building can be semantically obvious in pixels while remaining invisible to a
depth channel clipped at tens of meters.  Open-vocabulary models make this
evidence addressable in free language
\cite{clip,lseg,openseg,pearl}, and recent navigation systems already exploit
it to choose \emph{where to go} \cite{vlfm,rayfronts}.
What they largely do not address is \emph{where not to go} at flight speed:
their semantic layers run at a few hertz, hold no executable notion of free
space, and hand a reactive executor a goal rather than a route.  For an aerial
vehicle, the resulting gap between a slow semantic opinion and a fast
commitment is precisely where missions are lost.

ScaleNav closes this gap by making the hierarchy explicit rather than
implicit in a monolithic network.  Its first layer converts generic depth and
odometry into a persistent graph of verified free-space bubbles connected by
collision-checked witness polylines, with no occupancy grid or ESDF in the
route layer \cite{bubble,epic}.  The retained backbone may grow as new space is
explored, but rolling A* and graph updates operate on a fixed-radius local
working set.  Its second, slower layer performs language-conditioned routing
on that graph.  Depth-clipped semantic observations become provisional
far-field frontiers and a continuous risk field along entire witnesses.  A
topology-graph GCN ranks five frontier directions from the online skeleton,
while A*,
collision checks, and route memory remain authoritative.  The selected route
is stored as a world-frame witness and remapped after atomic graph replacement.
Its third, fast layer is deliberately replaceable: ScaleNav samples a moving
local goal along the accepted witness and publishes a standard world-frame
goal message.  YOPO \cite{yopo} consumes that goal directly; EGO-Planner and
SUPER consume it through thin topic adapters.  Their planner cores are neither
retrained nor made graph-aware.

This separation lets us state the safety boundary precisely, and we regard
that boundary as a first-class design contribution.  The mapping layer
validates measured free space.  A provisional semantic-frontier connection is
an optimistic route hypothesis: it is checked against all currently known
obstacles before publication, but it is never treated as proof that
unobserved space is free.  Semantic risk repels or attracts route choices,
while current depth and the local executor remain solely responsible for
newly observed obstacles.  In the language of risk-aware planning
\cite{blackmore}, semantic evidence enters as a shaped preference
over verified alternatives, not as a probabilistic safety guarantee.

The main contributions are:
\begin{itemize}
  \setlength\itemsep{0pt}
  \item \textbf{A plug-and-play persistent topological route layer.}
  Verified bubbles, witness edges, stable identities, and atomic replacement
  preserve route structure.  A moving world-frame local goal exposes this
  memory to goal-conditioned planners without graph inputs or retraining,
  while fixed-radius operations control online work
  (Secs.~\ref{sec:topology} and~\ref{sec:interface}).
  \item \textbf{A topology-graph GCN frontier policy.}
  Weighted graph convolution, global topology pooling, and temporal direction
  state learn a five-way, 35-m-look-ahead frontier decision directly on the
  online skeleton, while downstream A* retains reachability and collision
  authority (Sec.~\ref{sec:gcn}).
  \item \textbf{Language preference separated from geometric authority.}
  Provisional far-field frontiers and witness-level semantic risk can change
  route ranking before reliable depth, while A* connectivity and collision
  checks retain control of route acceptance (Sec.~\ref{sec:semantics}).
  \item \textbf{Cross-executor validation.}
  The same topology-layer implementation is connected to YOPO, EGO-Planner,
  and SUPER through executor-specific goal and command adapters, and evaluated
  with conventional closed-loop flight and online-resource metrics
  (Sec.~\ref{sec:experiments}).
\end{itemize}

The evaluation separates system integration, route selection, and graph
persistence.  All reported semantic observations use PEARL \cite{pearl}; the
interface itself is model-agnostic.

%%%%%%%%%%%%%%%%%%%%%%%%%%%%%%%%%%%%%%%%%%%%%%%%%%%%%%%%%%%%%%%%%%%%%%%%%%%%%%%%
\section{RELATED WORK}
Reactive aerial planners trade global consistency for latency: RAPPIDS and
NanoMap operate on depth and short history \cite{rappids,nanomap}, EGO-Planner
removes the ESDF from optimization \cite{egoplanner}, and learned policies
generate motion primitives directly \cite{yopo,loquercio}.  Fast-Planner,
FASTER, and SUPER instead combine search or point-cloud corridors with
trajectory refinement \cite{fastplanner,faster,super}.  ScaleNav is
orthogonal: it supplies a persistent route decision above an unchanged local
planner core through a goal adapter.

Topological and frontier memories provide the natural substrate for this
separation.  Incremental frontier structures and bubble graphs
\cite{yamauchi,tare,epic,bubble} preserve reusable free-space structure,
while learned topological memories \cite{sptm,ving} show the value
of graph state for long-horizon navigation.  Our graph stores verified
witnesses, preserves route identity across atomic replacement, and exposes
the same skeleton topology to the GCN frontier policy.

Graph convolutional networks aggregate relational state directly over sparse
connectivity \cite{kipf2017gcn}, making them a natural alternative to raster or
independent-candidate encoders for route selection.  ScaleNav applies weighted
graph convolution to the online navigation skeleton, augments graph context
with temporal direction state, and learns a five-way frontier policy while
leaving route construction and collision validation outside the network.

Open-vocabulary models align image regions with text
\cite{clip,lseg,openseg,pearl} and have been lifted into semantic maps.
Frontier-selection systems \cite{vlfm,rayfronts} use such evidence to choose
where to go.
ScaleNav addresses the complementary aerial problem of where to avoid:
semantic heatmaps become provisional graph frontiers and witness-level risk,
but geometric graph checks remain authoritative.

Risk-aware planning shapes route preference without treating confidence as a
collision probability \cite{blackmore}.  We follow this
principle and keep semantic evidence behind a narrow interface; the
topology graph, A*, and collision checks define executable free space.
%%%%%%%%%%%%%%%%%%%%%%%%%%%%%%%%%%%%%%%%%%%%%%%%%%%%%%%%%%%%%%%%%%%%%%%%%%%%%%%%
\section{METHOD}
\label{sec:method}
\subsection{Problem Formulation and System Overview}
Let the vehicle state at time $t$ be
$\vect{x}_t=(\vect{p}_t,\vect{v}_t,\vect{a}_t,\vect{R}_t)$.
The sensors provide an RGB image $I_t$, a depth image $D_t$, and odometry.
The mission specifies a world-frame goal $\vect{g}$ and a text query $q$
describing regions to avoid.  No prior geometric map is available.

The planner must output a dynamically feasible local trajectory $\tau_t$.
Following the layered organization, mapping, route selection, and execution
compose as
\begin{align}
  \mathcal{G}_t
    &= \Pi_{\mathrm{map}}(\mathcal{G}_{t-1},D_t,\vect{x}_t), \\
  (\mathcal{W}_t,\vect{g}^{\mathrm{front}}_t,
   \vect{g}^{\mathrm{loc}}_t,\hat c_t)
    &= \Pi_{\mathrm{route}}(\mathcal{G}_t,I_{0:t},q,\vect{g}), \\
  \tau_t
    &= \Pi_{\mathrm{exec}}(D_t,\vect{o}_t).
\end{align}
Each layer has a deliberately narrow interface.  $\Pi_{\mathrm{map}}$ consumes
depth, odometry, and its previous graph state; neither the semantic model nor
the local planner appears in its update.  The slower $\Pi_{\mathrm{route}}$
produces an accepted world-frame witness $\mathcal{W}_t$, an internal route
anchor $\vect{g}^{\mathrm{front}}_t$, and a moving execution goal
$\vect{g}^{\mathrm{loc}}_t$ sampled ahead on that witness.  When enabled, the
topology-graph policy additionally publishes a five-way frontier column
$\hat c_t\in\{0,1,2,3,4\}$; this is a ranking hint, not a trajectory or safety
decision.  The plug-and-play boundary is the single standard goal
$\vect{g}^{\mathrm{loc}}_t$.  For YOPO, the released policy receives its usual
$160\times96$ depth channel, body-frame motion, and this goal expressed in the
body frame:
\begin{equation}
 \vect{o}_t = [\vect{v}^{B}_t,\vect{a}^{B}_t,
                (\vect{g}^{\mathrm{loc}}_t-\vect{p}_t)^{B}] \in \R^9.
\end{equation}
EGO-Planner and SUPER receive the same world-frame goal through a thin ROS
topic adapter.  This division mirrors hierarchical topological control
\cite{ving}, but no graph tensor, semantic feature, witness polyline, or
corridor representation crosses the executor boundary.

\begin{figure*}[!t]
\centering
\safefig[width=0.98\textwidth]{pics/system_architecture_topology_v3.pdf}
\vspace{-2pt}
\caption{ScaleNav as a plug-and-play route layer.  Depth and odometry maintain
persistent verified topology; RGB and text shape route cost, and a topology
GCN ranks frontier directions.  A* and collision checks accept a
witness before the goal interface drives YOPO, EGO-Planner, or SUPER.  Only
simulator I/O adapters are executor-specific.}
\label{fig:system_architecture}
\end{figure*}

We use a fixed-height planning layer in the current experiments to isolate
the horizontal branch-selection problem.  Depth sensing, point memory, and
local trajectory execution remain 3-D, and the semantic cost is defined on
3-D witness polylines without requiring the fixed-height assumption.

%%%%%%%%%%%%%%%%%%%%%%%%%%%%%%%%%%%%%%%%%%%%%%%%%%%%%%%%%%%%%%%%%%%%%%%%%%%%%%%%
\subsection{Topology Graph Construction and Maintenance}
\label{sec:topology}
The mapper receives depth and odometry and publishes an atomic skeleton graph.
It is independent of the semantic model, GCN, and local executor.  The current
depth frame is voxel-filtered into obstacle points and free-ray evidence; no
occupancy grid or ESDF is retained.  Collision-free bubbles are clustered into
vertices
\begin{equation}
 v_i=(\vect{c}_i,\rho_i,z_i,h_i),
\end{equation}
where $\vect{c}_i$ is center, $\rho_i$ is clearance, $z_i$ is semantic state,
and $h_i$ is a persistent identity.  Each edge stores a checked witness
polyline $\pi_{ij}$ and length $L_{ij}$.  The witness, rather than the chord
between endpoints, is the geometric object used by both collision checking
and later semantic fusion.

Graph construction runs asynchronously.  A background worker builds bubbles,
clusters regions, connects candidate edges, and atomically swaps a complete
snapshot into the planner.  The point buffer, local radius, and virtual-node
cap bound transient memory; persistent identities survive graph replacement,
while stale bubbles receive a grace period and are then deleted.  The minimum
witness clearance $\bar\rho_{ij}$ contributes a soft preference
\begin{equation}
 C_{\mathrm{clr}}(e_{ij})=
 \lambda_{\mathrm{clr}}L_{ij}
 \left(\frac{\rho^\star}{\rho^\star+\bar\rho_{ij}}\right)^2,
\label{eq:clearance}
\end{equation}
but never turns a geometrically valid edge into an invalid one.
%%%%%%%%%%%%%%%%%%%%%%%%%%%%%%%%%%%%%%%%%%%%%%%%%%%%%%%%%%%%%%%%%%%%%%%%%%%%%%%%
\subsection{Semantic Heatmap--Topology Fusion}
\label{sec:semantics}
The frozen PEARL model maps $(I_t,q)$ to a patch heatmap.  We max-pool it on
a $5\times3$ grid and project each retained patch at fixed optical-axis depth
$\ell_s$:
\begin{equation}
 \vect{y}_{u,t}=\vect{p}_t+\mat{R}_t
 \left(\vect{t}_{\mathrm{cam}}+\ell_s\widetilde{\vect{d}}_u\right).
\label{eq:projection}
\end{equation}
The projection does not consult measured depth, so an empty far-field return
cannot erase semantic evidence.  A calibrated score and confidence tuple
$(s_i,c_i)$ is associated with a persistent graph identity; stale observations
are rejected by timestamp gates and graph replacement restores the tuple.

High-scoring, spatially separated projections become provisional semantic
frontiers.  An unknown frontier may attach to at most two nearby verified
backbone nodes, but its witness is checked against all currently known
obstacles and cannot connect two verified corridors through unknown space.
When later geometry creates a bubble at the same position, the identity is
promoted to a verified node.  Thus a frontier is a route hypothesis, never a
free-space certificate.

Semantic evidence is fused with complete witness edges rather than only with
endpoints.  For semantic node $k$ and edge witness $\pi_{ij}$,
\begin{equation}
\begin{aligned}
 d_{k,ij}&=\min_{\vect{x}\in\pi_{ij}}
             \|\vect{y}_k-\vect{x}\|_2,\\
 r_{ij}&=\max\Big\{\tfrac12(s_i c_i+s_jc_j),\\[-2pt]
 &\qquad \max_{k:d_{k,ij}\le R_s}s_kc_k
 \exp[-d_{k,ij}^2/(2(R_s/2)^2)]\Big\}.
\end{aligned}
\label{eq:semantic_fusion}
\end{equation}
The resulting preference cost is
\begin{equation}
 C_{\mathrm{sem}}(e_{ij})=
 \lambda_{\mathrm{sem}}L_{ij}
 [-\log(\max(\epsilon,1-r_{ij}))].
\label{eq:semantic_cost}
\end{equation}
It reshapes route preference but never overrides geometric validity.

The planner searches reachable graph vertices within radius $R_g$ using
$C(e)=W_{ij}+C_{\mathrm{clr}}(e)+C_{\mathrm{sem}}(e)$.  It stores the selected
witness in world coordinates and holds it until progress reaches a fixed
fraction or its geometric check fails.  A fresh search then replaces the
complete route; no graph-node pointers are retained across atomic replacement.
The asynchronous contract uses timestamped snapshots, so graph, semantic, and
control streams cannot expose a partially rebuilt topology.
%%%%%%%%%%%%%%%%%%%%%%%%%%%%%%%%%%%%%%%%%%%%%%%%%%%%%%%%%%%%%%%%%%%%%%%%%%%%%%%%
\subsection{Planner-Agnostic Goal Interface}
\label{sec:interface}
The accepted witness, rather than a modified local policy, carries the route
decision.  Let $\mathcal{W}_t(s)$ be its arc-length parameterization and
$s_t$ the closest forward progress to the vehicle.  ScaleNav publishes
\begin{equation}
 \vect{g}^{\mathrm{loc}}_t=
 \mathcal{W}_t\!\left(\min(s_t+\ell_{\mathrm{loc}},L_{\mathcal W})\right),
 \qquad \ell_{\mathrm{loc}}=15~\mathrm{m}.
\label{eq:local_goal}
\end{equation}
The farther $\vect{g}^{\mathrm{front}}_t$ identifies the selected route branch
inside ScaleNav and is published for state inspection; the executor follows
only $\vect{g}^{\mathrm{loc}}_t$.  When a witness is replaced, the vehicle is
projected onto the new polyline before Eq.~\eqref{eq:local_goal} is evaluated.
A short hold gate suppresses transient goal loss during an atomic graph swap.

The implementation exposes $\vect{g}^{\mathrm{loc}}_t$ as a world-frame
\texttt{geometry\_msgs/PoseStamped}.  The original YOPO checkpoint consumes it
through its ordinary relative-goal input.  For EGO-Planner and SUPER, a thin
bridge rate-limits and republishes the same message on their native goal topic;
separate point-cloud and command bridges adapt the shared simulator I/O.  These
adapters do not alter either planner's search, optimization, or control law.
Consequently, the transferable contract is precise: a local planner must
accept a position goal and publish executable commands; it need not understand
ScaleNav's graph, semantics, or route memory.

The closed-loop experiments pass no corridor, bubble sequence, route mask, or
graph feature into any planner and do not retrain planner weights.  ScaleNav
instead preserves route intent by advancing the goal along the witness, while
current depth and the local planner remain responsible for short-horizon
motion.  A stale or invalid witness is rejected upstream and triggers new A*.

%%%%%%%%%%%%%%%%%%%%%%%%%%%%%%%%%%%%%%%%%%%%%%%%%%%%%%%%%%%%%%%%%%%%%%%%%%%%%%%%
\subsection{Topology-Graph GCN Frontier Policy}
\label{sec:gcn}
The GCN is ScaleNav's learned long-range route-selection component.  It
predicts which of five directions should be considered first: left $40^\circ$,
left $20^\circ$, forward, right $20^\circ$, or right $40^\circ$.  Its output is
deliberately constrained to a frontier proposal: it neither creates a route
nor publishes an executor command, leaving geometric acceptance to A*.

At each inference tick, skeleton nodes and edges form the input graph.  Five
virtual candidates are inserted $10$~m from the vehicle and connected to their
nearest skeleton node with inverse-distance edge weights.  Each node has a
20-dimensional feature vector containing normalized world and body-frame
position; bubble radius and degree; semantic score and confidence; candidate
identity; origin and goal distance; odometry and forward-candidate flags; yaw;
range; bearing; and body-frame mission-goal state.  The deployed
\texttt{FrontierGCN(input\_dim=20)} has two weighted GCN layers, global mean
pooling, a GRU over the previous direction, and an MLP candidate scorer
\cite{kipf2017gcn}.  Its
column is published on \texttt{/scalenav/gcn\_frontier\_column}; A* then tests
reachability and collision validity exactly as it does for a hand-designed
ranking.

Training labels use the complete static Map2 point cloud, projected onto a
$0.75$~m grid and inflated by the $1.2$~m vehicle radius.  Eight-neighbor A*
provides the direction of the first path point at least $35$~m ahead;
out-of-view and incomplete paths are discarded.  Session-level splits and
class-weighted cross-entropy prevent adjacent frames from crossing splits.
The resulting classifier is evaluated separately in
Sec.~\ref{sec:heuristic_eval}; its accuracy is not a safety claim.

%%%%%%%%%%%%%%%%%%%%%%%%%%%%%%%%%%%%%%%%%%%%%%%%%%%%%%%%%%%%%%%%%%%%%%%%%%%%%%%%
\section{EXPERIMENTS}
\label{sec:experiments}
The evaluation examines three complementary aspects of the system: closed-loop
compatibility of the topological goal interface with different local planners,
the effects of semantic cost and the GCN policy on route selection,
and the online graph-search workload under persistent geometry.
Closed-loop and profiling trials use UE Map2 through Colosseum/AirSim
\cite{airsim} with synchronized RGB-D and odometry and a 1.6-m planning
layer.  One-way trials fly from $(0,0,1.6)$ to $(0,140,1.6)$~m.

\subsection{Evaluation Protocol}
\label{sec:protocol}
\emph{Statistics.}  The reported batches were collected separately and are
not seed matched, so all cross-condition comparisons are descriptive.  Startup
failures are excluded from navigation outcomes, while collision and timeout
remain in the outcome rates; completion metrics use successful trials only.

\emph{Metrics.}  Closed-loop metrics are success, collision, timeout, path
length $L$, flight time $T$, $\bar v=L/T$, and maximum speed.  For failed
flights, distance traveled before termination is reported separately as
$L_{\rm obs}$ and never mixed into successful-flight completion means.
Profiling separates point processing, graph construction, active A*, and the
complete planning tick.

\emph{Baseline fairness.}  EGO-Planner \cite{egoplanner} and SUPER
\cite{super} retain their planner cores but require simulator-specific cloud,
goal, and command adapters.  Their Map2 configurations and timeouts differ
from YOPO's, and recent benchmarking shows that such stacks are sensitive to
latency, field of view, and scenario difficulty \cite{flightbench}.  Their
failures therefore characterize these deployments, not universal planner
capability.  The plug-in claim concerns the narrow goal interface and absence
of planner-core changes, not strict experimental interchangeability.

\subsection{Cross-Planner Closed-Loop Evaluation}
\subsubsection{Methods and Conditions}
The primary comparison pairs each standalone executor with a ScaleNav-guided
variant.  All ScaleNav variants run the same graph-node implementation,
semantic route formulation, A* acceptance, witness memory, and local-goal
publisher.  YOPO consumes the goal directly; EGO-Planner and SUPER use the
adapter described in Sec.~\ref{sec:interface}.  Their internal planners do not
consume ScaleNav data structures.  Each row in the plug-in comparison contains
ten valid trials.  Table~\ref{tab:params} lists route-layer parameters shared
by the principal YOPO ablations.

\begin{table}[t]
\caption{Fixed Parameters Used by the Current Implementation}
\label{tab:params}
\centering
\small
\resizebox{\columnwidth}{!}{%
\begin{tabular}{lc@{\quad}lc}
\toprule
Parameter & Value & Parameter & Value \\
\midrule
Point voxel & 0.10 m & Point budget & 100,000 \\
Frame history & current frame only & Graph update & 100 ms \\
Graph rebuild & 100 ms & Search radius & 45 m \\
Local lookahead & 15 m & Goal-connect budget & 20 ms \\
$\lambda_{\rm sem}$ & 2.0 & $R_s$ & 5 m \\
$\lambda_{\rm clr}$ & 2.0 & $\rho^\star$ & 1.2 m \\
Score update & latest match & Edge-memory discount & disabled \\
Optical depth $\ell_s$ & 30 m & Node merge distance & 2.5 m \\
Semantic candidates/frame & 16 & Virtual-node cap & 512 \\
Semantic candidates/edge & 8 & Semantic grid & $5\times3$ \\
\bottomrule
\end{tabular}
}
\end{table}

\subsubsection{Scenario}
Map2 contains isolated and dense trees, long walls, buildings, and
mixed-scale clusters (Fig.~1(b)).  Its truth export spans
$140.0\times132.0$~m.  In the evaluation corridor $x\in[-45,45]$~m,
$y\in[0,120]$~m, $3.91\%$ of 0.5-m cells intersect truth geometry in the
navigation slab $z\in[0.6,2.6]$~m.  These statistics come from UE truth,
never from any method's partial online map.  Matched language trials
use the query ``blocks, wall.''

\subsubsection{Results}
Figure~\ref{fig:map2_speed_trajectories} is a qualitative visualization of
representative current-scene trajectories.  It contrasts frontier-only runs
with graph-guided ScaleNav runs for YOPO, EGO-Planner, and SUPER; the
ScaleNav (YOPO) panel uses the shortest successful trajectory from the latest
ten-trial batch.  Aggregate results are reported in
Table~\ref{tab:graph_collision_comparison}.

\begin{figure*}[!ht]
\centering
\safefig[width=0.90\textwidth]{pics/experiments/map2_0_140_1p6/trajectory_speed_comparison_current.pdf}
\vspace{-2pt}
\caption{Representative Map2 trajectories with and without the topology graph.
Rows compare YOPO-Simple,
EGO-Planner, and SUPER; the left column uses each standalone planner with its
direct mission goal and the right column uses the ScaleNav route layer.  The
ScaleNav (YOPO) trajectory is the shortest successful run from the latest
ten-trial GCN batch, and the ScaleNav (EGO) trajectory is the shortest
successful run from the latest ten-trial regression.  Color denotes executed speed; circles,
stars, and crosses mark starts, completed goals, and incomplete collision
endpoints.}
\label{fig:map2_speed_trajectories}
\end{figure*}

\begin{table*}[t]
\caption{Plug-in evaluation on the Map2 long-range mission ($n=10$ per row).
Each pair contrasts the executor's standalone goal with ScaleNav's moving
local goal.  ScaleNav rows share the same topology-layer implementation;
executor-specific adapters only relay goals, point clouds, and commands.}
\label{tab:graph_collision_comparison}
\centering
\small
\resizebox{0.98\textwidth}{!}{%
\begin{tabular}{lccccccc}
\toprule
Method & Success (\%) & Collision (\%) & Timeout (\%) & Path (m) & Time (s) & Avg. speed (m/s) & Max. speed (m/s) \\
\midrule
YOPO-Simple (no graph) & 0 & 100 & 0 & $129.79\pm0.58^\dagger$ & $61.46\pm0.36^\dagger$ & $2.11\pm0.01^\dagger$ & $6.46\pm0.02^\dagger$ \\
ScaleNav (YOPO) & 100 & 0 & 0 & $171.59\pm6.57$ & $33.61\pm1.25$ & $5.105\pm0.045$ & $6.254\pm0.138$ \\
EGO-Planner (no graph) & 0 & 100 & 0 & $73.24\pm0.69^\dagger$ & $54.32\pm1.05^\dagger$ & $1.35\pm0.02^\dagger$ & $2.43\pm0.25^\dagger$ \\
ScaleNav (EGO) & 70 & 30 & 0 & $181.59\pm11.33$ & $114.05\pm7.82$ & $1.593\pm0.017$ & $2.852\pm0.292$ \\
SUPER (no graph) & 0 & 100 & 0 & $47.07\pm28.27^\dagger$ & $10.25\pm4.99^\dagger$ & $4.30\pm0.68^\dagger$ & $5.60\pm0.23^\dagger$ \\
ScaleNav (SUPER) & 80 & 10 & 10 & $181.40\pm16.96$ & $40.29\pm3.48$ & $4.502\pm0.140$ & $5.821\pm0.026$ \\
\bottomrule
\end{tabular}
}
\end{table*}

All three standalone batches terminated without a successful mission, whereas
the ScaleNav-guided YOPO, EGO-Planner, and SUPER variants completed 10/10,
7/10, and 8/10 trials.  This demonstrates that the same topological goal
contract is executable by three different planner cores.  It does not isolate
which route-layer component causes the outcome gap because the batches are
unmatched and the ScaleNav variants include semantic cost and the GCN ranking
heuristic.  Values marked $\dagger$ are means over truncated failed flights,
not completion metrics; unmarked continuous values use successful flights
only.  The ScaleNav (EGO) deployment uses a 200-s mission timeout, while the
other rows use 90~s.

The ScaleNav (YOPO) run produced no collision or timeout events: all ten trials
reached the goal and satisfied the final position and settling-speed
tolerances.  Relative to the 140-m straight-line mission, its mean path is
31.59~m longer, corresponding to a $22.6\%$ path overhead.  The longest trial
was trial 6 at $188.43$~m and $36.72$~s, which accounts for most of the path
length spread.  The mean terminal position error was $0.096\pm0.029$~m and
the mean terminal speed was $0.294\pm0.004$~m/s.

\subsection{Route-Layer Ablations}
\subsubsection{Semantic Influence}
To test whether the route benefit depends on semantic evidence beyond the
depth envelope, we repeated the GCN stack with the same Map2 mission,
90-s timeout, graph configuration, and query, while disabling PEARL,
semantic nodes, and semantic edge cost.  The two batches were collected
separately rather than seed-paired; the comparison is therefore an ablation
of the semantic interface, not a causal estimate over identical stochastic
trials.  Completion metrics use successful flights only, while $L_{\rm obs}$
records the path observed before collision or timeout.  The resulting
conventional flight metrics are reported in Table~\ref{tab:gcn_semantic_ablation}.

\begin{table}[t]
\caption{Semantic route-layer ablation on Map2 with the GCN heuristic fixed.
Semantic-off removes PEARL heatmaps, far-field nodes, and semantic edge cost;
geometry, A*, goal publication, and YOPO remain enabled.}
\label{tab:gcn_semantic_ablation}
\centering
\small
\resizebox{\columnwidth}{!}{%
\begin{tabular}{lcccccccc}
\toprule
Condition & Success (\%) & Collision (\%) & Timeout (\%) & $L_{\rm succ}$ (m) & $L_{\rm obs}$ (m) & Time (s) & Avg. speed (m/s) & Max. speed (m/s) \\
\midrule
Matched query & 100 & 0 & 0 & $171.59\pm6.57$ & \blank & $33.61\pm1.25$ & $5.105\pm0.045$ & $6.254\pm0.138$ \\
Semantic disabled & 60 & 40 & 0 & $169.47\pm3.16$ & $11.69\pm1.60$ & $32.34\pm2.70$ & $5.261\pm0.315$ & $6.404\pm0.158$ \\
Irrelevant query & 80 & 20 & 0 & $175.76\pm6.09$ & $20.32\pm2.67$ & $35.02\pm1.65$ & $5.023\pm0.127$ & $6.398\pm0.274$ \\
\bottomrule
\end{tabular}
}
\end{table}

Semantic influence changes the outcome rate without changing the geometric
validity contract: the semantic-on batch completed all ten flights, whereas
the semantic-off batch had four early collisions near $y\approx10$~m.  The
successful-flight path and speed differences are secondary to this outcome
gap and should not be read as a speed improvement.  Together with the
depth-clipped projection visualized in Fig.~1, this behavior is consistent
with semantic evidence affecting a route before the corresponding obstacle is
measured by depth; the unmatched trials do not establish a causal effect.
The irrelevant-query condition (prompt ``trees'') completes 8/10 trials,
showing that semantic performance also depends on query alignment with the
scene and should not be summarized as a query-independent guarantee.

\subsubsection{GCN Frontier-Policy Ablation}
To isolate the learned route-selection component, we compare ScaleNav's
hand-designed frontier ranking with the GCN policy while
holding the Map2 mission, semantic query, route-cost parameters, and YOPO
executor fixed.  The batches are unmatched and have different sample counts.

\begin{table}[t]
\caption{Frontier-Policy Ablation Within ScaleNav}
\label{tab:single_run_validation}
\centering
\small
\resizebox{\columnwidth}{!}{%
\begin{tabular}{lccccccc}
\toprule
Method & Success (\%) & Collision (\%) & Timeout (\%) & Path (m) & Time (s) & Avg. speed (m/s) & Max. speed (m/s) \\
\midrule
Hand-designed ($n=27$) & 88.9 & 11.1 & 0 & $199.14\pm14.12$ & $38.54\pm3.06$ & $5.172\pm0.097$ & $6.409\pm0.184$ \\
GCN ($n=10$) & 100.0 & 0.0 & 0.0 & $171.59\pm6.57$ & $33.61\pm1.25$ & $5.105\pm0.045$ & $6.254\pm0.138$ \\
\bottomrule
\end{tabular}
}
\end{table}

The GCN batch completed 10/10 trials, compared with 24/27 for the hand-designed
ranking.  Its shorter mean path and time are descriptive rather than causal;
the comparison shows the effect of replacing the hand-designed ranking with
the GCN policy while A*, the accepted witness, and the local-goal interface
remain unchanged.

\subsection{GCN Policy Evaluation}
\label{sec:heuristic_eval}
The privileged Map2 dataset contains 5,337 training samples and 707 held-out
test frames.  The plain checkpoint obtains $89.3\%$ exact and $89.4\%$ macro
accuracy, while validation-selected switch penalty $0.05$ gives $91.3\%$ exact
and $91.8\%$ macro accuracy.  The original planner matches only $22.2\%$ of
the strict 35-m labels.  These results measure route-direction prediction,
not online safety.

\begin{table}[t]
\caption{GCN Direction Classification}
\label{tab:gcn_offline}
\centering
\small
\begin{tabular}{lcc}
\toprule
Metric & Plain & Penalty $0.05$ \\
\midrule
Test exact accuracy & 89.3\% & 91.3\% \\
Test macro accuracy & 89.4\% & 91.8\% \\
Within one direction column & 99.4\% & \blank \\
\bottomrule
\end{tabular}
\end{table}

Figure~\ref{fig:gcn_real_views} shows synchronized frames from the online
interface.  The RGB image, PEARL heatmap, selected GCN column, and mission-goal
direction are retained from the recorded Map2 experiment.

\begin{figure*}[t]
\centering
\begin{minipage}[t]{0.235\textwidth}
\centering
\safefig[width=\linewidth]{pics/0903_graph_semantic_frontier_building1.png}\\[-1pt]
{\footnotesize (a) Building view 1}
\end{minipage}\hfill
\begin{minipage}[t]{0.235\textwidth}
\centering
\safefig[width=\linewidth]{pics/0903_graph_semantic_frontier_building2.png}\\[-1pt]
{\footnotesize (b) Building view 2}
\end{minipage}\hfill
\begin{minipage}[t]{0.235\textwidth}
\centering
\safefig[width=\linewidth]{pics/0903_graph_semantic_frontier_forest1.png}\\[-1pt]
{\footnotesize (c) Forest view 1}
\end{minipage}\hfill
\begin{minipage}[t]{0.235\textwidth}
\centering
\safefig[width=\linewidth]{pics/0903_graph_semantic_frontier_forest2.png}\\[-1pt]
{\footnotesize (d) Forest view 2}
\end{minipage}
\caption{Recorded GCN-interface views from Map2 flights.  Each panel preserves
the RGB frame and PEARL heatmap; red marks the proposed frontier column and the
green dashed arrow marks the 35-m mission-goal direction.}
\label{fig:gcn_real_views}
\end{figure*}

\subsection{Online Planning Workload}
The online search is restricted to a local working set even when the verified
topology persists outside the current sensing window.  The current
configuration keeps at most 100,000 point samples, rebuilds the skeleton every
100~ms, searches within 45~m, and caps virtual semantic nodes at 512.
Figure~\ref{fig:graph_resource_scaling} compares the persistent-geometry mode
with a sliding-geometry ablation.  The persistent mode uses 214/249 local
nodes at P50/P95 and expands 185/215 nodes per active A* search, compared with
197/231 and 173/194 when persistence is disabled.  Thus retaining the explored
backbone does not make every retained node part of each planning tick.  Over
normalized session time, the observed local-node P95 remains below 279 with
persistence and 248 without it.

\begin{figure}[t]
\centering
\safefig[width=\columnwidth]{pics/experiments/map2_0_140_1p6/graph_resource_scaling.pdf}
\vspace{-3pt}
\caption{Online resources for persistent and 40-m sliding geometry.
(a)~Local nodes over normalized time. (b)~Module wall time.  Lines/bars show
P50/mean and dashed lines/caps show P95.  The unmatched profiles use 27 and 10
Map2 sessions; global nodes are excluded from the local search.}
\label{fig:graph_resource_scaling}
\end{figure}

A route holds its accepted witness until a progress gate is reached or a
geometric check fails; this makes route switching observable and prevents a
partially rebuilt graph from leaking into the controller.  The semantic stream
runs at roughly 2~Hz, whereas graph and control updates remain asynchronous.

For reproducibility, each trial records the graph snapshot timestamp, selected
frontier column, witness clearance, A* replacement reason, local-goal distance,
and terminal outcome.  These fields let us distinguish a direction proposal
from the route actually accepted after collision checking.  They also expose
the failure sequence seen in the current batch: a stale or low-clearance
witness is rejected, a fresh A* search is requested, and a moving vehicle may
reach the executor's start-failure condition before the replacement goal is
established.  This logging boundary is important because it attributes safety
events to the graph--planner--executor interface rather than to an offline GCN
classification score alone.

\subsection{Statistical Scope}
Each cross-executor row contains ten valid navigation trials, but paired rows
are not seed matched and the executor deployments do not share every timeout
or internal mapping parameter.  The heuristic comparison contains 27 valid
hand-designed-ranking trials and ten GCN trials; one manual interruption and
startup failures are excluded.  The semantic conditions contain ten trials
each but are likewise unmatched.  The results therefore demonstrate interface
compatibility and report observed outcome rates; they do not provide a paired
causal estimate of improvement.
%%%%%%%%%%%%%%%%%%%%%%%%%%%%%%%%%%%%%%%%%%%%%%%%%%%%%%%%%%%%%%%%%%%%%%%%%%%%%%%%
\section{DISCUSSION AND LIMITATIONS}
\label{sec:discussion}
The contribution is the topology route layer and its narrow goal contract, not
a modified local policy.  The current YOPO checkpoint receives only depth,
state, and the moving local goal.  EGO-Planner and SUPER require simulator and
topic adapters, but their planner cores do not receive graph features or
semantic scores.  Thus ``plug-and-play'' means no
planner retraining or algorithm modification after the goal and command
interfaces are adapted; it does not mean arbitrary planners can be connected
without integration code.

The single-goal interface is intentionally weak.  Unlike a corridor constraint,
it cannot force an executor to remain inside every bubble or inherit the
topological layer's collision checks.  ScaleNav instead advances the goal along
the accepted witness, while each executor retains responsibility for feasible
short-horizon motion.  The 10/10, 7/10, and 8/10 ScaleNav outcomes demonstrate
portability in Map2, not a common safety guarantee across executors.

The offline GCN accuracy does not by itself establish online safety.  The
latest ten-trial batch completed without collision, but it is not seed-matched
to the 27-trial original batch and is too small to support a significance
claim.  Future graph policies should still expose reachable clearance,
braking margin, and A* cost, and unsafe columns should be masked before
learned ranking.

The current evaluation is in one simulated Map2 environment, uses a
fixed-height planning layer, and compares unmatched batches.  Semantic
observations are evaluated with PEARL at roughly 2 Hz, and provisional
far-field nodes never certify unknown free space.  The retained global
backbone has no hard node limit; only the point buffer, semantic hypotheses,
and local working-set radius have explicit limits; active nodes are profiled.
Longer missions,
true process RSS, equal-sized seed-matched trials, and cross-map evaluation are
needed before claiming asymptotic memory bounds or statistically reliable
navigation improvement.

%%%%%%%%%%%%%%%%%%%%%%%%%%%%%%%%%%%%%%%%%%%%%%%%%%%%%%%%%%%%%%%%%%%%%%%%%%%%%%%%
\section{CONCLUSION}
ScaleNav is a topological route layer rather than another local planner.  It
turns depth and odometry into persistent witness-verified structure, combines
language risk with a topology-graph GCN frontier policy upstream, and exports only
a moving world-frame local goal.  The original YOPO policy consumes that goal
directly, while goal adapters connect the same topology implementation to
EGO-Planner and SUPER without planner retraining.  In the tested Map2
deployments, the standalone rows record 0/10 completions, while the
corresponding ScaleNav-guided rows record 10/10, 7/10, and 8/10 and the
observed local search remains a few hundred nodes.  The
unmatched trials do not establish universal superiority, but they support the
central systems claim: long-horizon memory and language-conditioned route
choice can be added above interchangeable short-horizon executors.

{\scriptsize
\begin{thebibliography}{99}

\bibitem{yopo}
J. Lu, X. Zhang, H. Shen, L. Xu, and B. Tian,
``You Only Plan Once: A Learning-Based One-Stage Planner With Guidance
Learning,'' \emph{IEEE Robot. Autom. Lett.}, vol. 9, no. 7,
pp. 6083--6090, 2024.

\bibitem{rappids}
N. Bucki, J. Lee, and M. W. Mueller,
``Rectangular Pyramid Partitioning Using Integrated Depth Sensors (RAPPIDS):
A Fast Planner for Multicopter Navigation,'' \emph{IEEE Robot. Autom. Lett.},
vol. 5, no. 3, pp. 4626--4633, 2020.

\bibitem{nanomap}
P. R. Florence, J. Carter, J. Ware, and R. Tedrake,
``NanoMap: Fast, Uncertainty-Aware Proximity Queries With Lazy Search of
Local 3-D Data,'' in \emph{Proc. IEEE Int. Conf. Robot. Autom.}, 2018,
pp. 7631--7638.

\bibitem{egoplanner}
X. Zhou, Z. Wang, H. Ye, C. Xu, and F. Gao,
``EGO-Planner: An ESDF-Free Gradient-Based Local Planner for Quadrotors,''
\emph{IEEE Robot. Autom. Lett.}, vol. 6, no. 2, pp. 478--485, 2021.

\bibitem{loquercio}
A. Loquercio, E. Kaufmann, R. Ranftl, M. M\"uller, V. Koltun, and
D. Scaramuzza,
``Learning High-Speed Flight in the Wild,''
\emph{Sci. Robot.}, vol. 6, no. 59, 2021.

\bibitem{bubble}
Y. Ren, F. Zhu, W. Liu, Z. Wang, Y. Lin, F. Gao, and F. Zhang,
``Bubble Planner: Planning High-Speed Smooth Quadrotor Trajectories Using
Receding Corridors,'' in \emph{Proc. IEEE/RSJ Int. Conf. Intell. Robots
Syst.}, 2022, pp. 6332--6339.

\bibitem{epic}
S. Geng, Z. Ning, F. Zhang, and B. Zhou,
``EPIC: A Lightweight LiDAR-Based AAV Exploration Framework for Large-Scale
Scenarios,'' \emph{IEEE Robot. Autom. Lett.}, vol. 10, no. 5,
pp. 5090--5097, 2025.

\bibitem{fastplanner}
X. Zhou, J. Gao, C. Wang, Z. Wang, H. Ye, and F. Gao,
``Fast-Planner: A Robust and Efficient Trajectory Planner for Autonomous
Quadrotor Flight,'' in \emph{Proc. IEEE/RSJ Int. Conf. Intell. Robots Syst.},
2019, pp. 8455--8461.

\bibitem{faster}
J. Tordesillas, B. T. Lopez, J. P. How, and L. Carlone,
``FASTER: A Robust Fast-Planner for Safe Autonomous Navigation in Unknown
Environments,'' \emph{IEEE Trans. Robot.}, vol. 38, no. 2, pp. 1112--1131,
2022.

\bibitem{super}
Y. Ren, F. Zhu, G. Lu, Y. Cai, L. Yin, F. Kong, J. Lin, N. Chen, and
F. Zhang,
``Safety-Assured High-Speed Navigation for MAVs,''
\emph{Sci. Robot.}, vol. 10, no. 98, p. eado6187, 2025.

\bibitem{yamauchi}
B. Yamauchi,
``A Frontier-Based Approach for Autonomous Exploration,'' in \emph{Proc.
IEEE Int. Symp. Comput. Intell. Robot. Autom.}, 1997, pp. 146--151.

\bibitem{tare}
C. Cao, H. Zhu, H. Choset, and J. Zhang,
``TARE: A Hierarchical Framework for Efficiently Exploring Complex 3D
Environments,'' in \emph{Proc. Robot.: Sci. Syst.}, 2021.

\bibitem{voxblox}
H. Oleynikova, Z. Taylor, M. Fehr, R. Siegwart, and J. Nieto,
``Voxblox: Incremental 3D Euclidean Signed Distance Fields for On-Board MAV
Planning,'' \emph{IEEE/RSJ Int. Conf. Intell. Robots Syst. (IROS)}, 2017.

\bibitem{sptm}
N. Savinov, A. Dosovitskiy, and V. Koltun,
``Semi-Parametric Topological Memory for Navigation,'' in \emph{Proc. Int.
Conf. Learn. Represent.}, 2018.

\bibitem{ving}
D. Shah, B. Eysenbach, G. Kahn, N. Rhinehart, and S. Levine,
``ViNG: Learning Open-World Navigation With Visual Goals,'' in \emph{Proc.
IEEE Int. Conf. Robot. Autom.}, 2021, pp. 13215--13222.

\bibitem{kipf2017gcn}
T. N. Kipf and M. Welling,
``Semi-Supervised Classification with Graph Convolutional Networks,''
in \emph{Proc. Int. Conf. Learn. Represent.}, 2017.

\bibitem{clip}
A. Radford et al.,
``Learning Transferable Visual Models From Natural Language Supervision,''
in \emph{Proc. Int. Conf. Mach. Learn.}, 2021, pp. 8748--8763.

\bibitem{lseg}
B. Li, K. Q. Weinberger, S. Belongie, V. Koltun, and R. Ranftl,
``Language-Driven Semantic Segmentation,'' in \emph{Proc. Int. Conf. Learn.
Represent.}, 2022.

\bibitem{openseg}
G. Ghiasi, X. Gu, Y. Cui, and T.-Y. Lin,
``Scaling Open-Vocabulary Image Segmentation With Image-Level Labels,'' in
\emph{Proc. Eur. Conf. Comput. Vis.}, 2022.

\bibitem{pearl}
G. Pei, X. Jiang, X. Cai, T. Chen, Y. Yao, and B. Jeon,
``PEARL: Geometry Aligns Semantics for Training-Free Open-Vocabulary Semantic
Segmentation,'' \emph{arXiv preprint arXiv:2603.21528}, 2026.

\bibitem{vlfm}
N. Yokoyama, S. Ha, D. Batra, J. Wang, and B. Bucher,
``VLFM: Vision-Language Frontier Maps for Zero-Shot Semantic Navigation,''
in \emph{Proc. IEEE Int. Conf. Robot. Autom.}, 2024, pp. 42--48.

\bibitem{rayfronts}
O. Alama, A. Bhattacharya, H. He, S. Kim, Y. Qiu, W. Wang, C. Ho, N. Keetha,
and S. Scherer,
``RayFronts: Open-Set Semantic Ray Frontiers for Online Scene Understanding
and Exploration,'' in \emph{Proc. IEEE/RSJ Int. Conf. Intell. Robots Syst.},
2025, pp. 5930--5937.

\bibitem{blackmore}
L. Blackmore, M. Ono, and B. C. Williams,
``Chance-Constrained Optimal Path Planning With Obstacles,'' \emph{IEEE
Trans. Robot.}, vol. 27, no. 6, pp. 1080--1094, 2011.

\bibitem{flightbench}
S.-A. Yu, C. Yu, F. Gao, Y. Wu, and Y. Wang,
``FlightBench: Benchmarking Learning-Based Methods for Ego-Vision-Based
Quadrotors Navigation,'' in \emph{Proc. IEEE Int. Conf. Robot. Autom.},
2025.

\bibitem{airsim}
S. Shah, D. Dey, C. Lovett, and A. Kapoor,
``AirSim: High-Fidelity Visual and Physical Simulation for Autonomous
Vehicles,'' in \emph{Field and Service Robotics}, 2018, pp. 621--635.

\end{thebibliography}
}

\end{document}
